%% Chapter 4: Demonstration Plan

\chapter{DEMONSTRATION PLAN}

This chapter describes how the proposed system will be demonstrated to the project panel.
The demonstration follows a two-phase structure: Phase~I demonstrates the complete
algorithmic pipeline through software simulations, and Phase~II
extends the validated software to hardware deployment and outdoor flight tests.

%% ─────────────────────────────────────────────────────────────────────────────
\section{PHASE I DEMONSTRATION --- SOFTWARE SIMULATION}

Phase~I demonstrations are entirely software-based. All demos run on a standard laptop
with no GPU or cloud dependency. The goal of Phase~I demos is to let the panel
\textbf{see the algorithm making decisions in real time} and to compare the proposed
system against baselines with quantitative output.

\subsection{Demo 1: Airspace Graph and Objective Visualisation}

\textbf{What is shown.} A live visualisation of the extracted OpenStreetMap airspace
graph for a chosen delivery region. Edge colours represent the three objective costs:
blue for distance, green for energy, red for risk. The operator can toggle between
objective overlays and zoom into specific delivery segments.

\textbf{Purpose.} Establishes that the Tier~1 inputs (the three-cost-vector graph) are
correctly constructed and visually interpretable. The panel can see how obstacle zones,
no-fly corridors, and wind-exposed segments translate into high-risk edge colours.

\textbf{Tools.} Python (\texttt{osmnx}, \texttt{networkx}, \texttt{matplotlib}).

\subsection{Demo 2: Tier 1 Pareto Front Generation and 3D Scatter Plot}

\textbf{What is shown.} The multi-objective A* algorithm is run on the medium-scale
scenario (S2, $\sim$200 nodes). The resulting Pareto front is visualised as an
interactive 3D scatter plot with axes: route distance, energy consumption, and flight
risk. All non-dominated routes are shown as points. The MOPSO baseline front is overlaid
in a contrasting colour for visual comparison.

\textbf{Purpose.} Demonstrates that Tier~1 produces a well-distributed Pareto front with
higher hypervolume (HV) than MOPSO. The panel can visually verify that no point in the
proposed system's front is dominated by any point in the MOPSO front.

\textbf{Tools.} Python (\texttt{matplotlib} 3D scatter, or \texttt{plotly} for
interactive rotation).

\subsection{Demo 3: Tier 2 Payload-Aware Route Selection (Interactive)}

\textbf{What is shown.} The ground-station dashboard interface is opened. The operator
enters payload tags for three scenarios in sequence:

\begin{enumerate}
  \item \textbf{Fragile medical sample}: $F=5, T=1, U=1, P=3$
  \item \textbf{Emergency medicine}: $F=1, T=1, U=3, P=5$
  \item \textbf{Disaster-relief supplies}: $F=1, T=2, U=1, P=3$
\end{enumerate}

For each tag entry, the system instantly:
\begin{itemize}
  \item Displays the derived weights ($w_{\mathrm{risk}}, w_{\mathrm{dist}}, w_{\mathrm{energy}}$)
  \item Highlights the selected Pareto point on the 3D scatter plot
  \item Overlays the selected route on the airspace graph map
  \item Shows the mission statistics: distance, estimated energy use, and risk score
\end{itemize}

\textbf{Purpose.} Provides the most direct visual demonstration of the core contribution.
The panel can see the selected point \textbf{move to a different, non-extreme location}
on the Pareto front for each payload type, confirming that the utility scoring selects
genuinely different trade-offs.

\textbf{Tools.} Tkinter or React dashboard; \texttt{matplotlib} or \texttt{plotly} for
real-time scatter highlight.

\subsection{Demo 4: Quantitative Benchmark Comparison}

\textbf{What is shown.} A table and bar chart comparing the proposed Tier~1 system
against single-objective A* and MOPSO on all three scenario scales (S1, S2, S3).
Metrics displayed: Hypervolume (HV), Inverted Generational Distance (IGD), and
computation time. MOPSO results include mean $\pm$ standard deviation over 30 runs.

\textbf{Purpose.} Provides the panel with rigorous quantitative evidence that Tier~1
produces a higher-quality Pareto front than both baselines, particularly on hypervolume
(which measures the volume of dominated objective space---a higher value is strictly
better).

\textbf{Tools.} Python benchmarking suite; \texttt{matplotlib} bar charts.

\subsection{Demo 5: Tier 2 Selection Accuracy Measurement}

\textbf{What is shown.} For each of the four payload presets (Table~\ref{tab:presets}),
the system's selected route is compared against the \textit{ground-truth optimal} route
(determined by exhaustive evaluation of all Pareto points with oracle weights). The
\textbf{Tier~2 selection accuracy} is reported as the percentage of payload types for
which the algorithm selected the same route as the oracle.

\textbf{Purpose.} Quantifies the alignment between payload-derived weights and the
routing requirements they express, establishing that the algebraic mapping is not just
principled in theory but correct in practice.

\textbf{Tools.} Python evaluation script; tabular output.

%% ─────────────────────────────────────────────────────────────────────────────
\section{PHASE II DEMONSTRATION --- HARDWARE IMPLEMENTATION}
\label{sec:phase2demo}

Phase~II demonstrations are conducted once the software system has been fully validated in Phase~I. 
All Phase~I algorithmic validation is completed before any hardware is powered on. 
Real flight tests confirm an already-validated system, not the primary evaluation environment.

\subsection{Step 1: ArduPilot SITL Validation}

Before outdoor flight, every payload preset is tested in ArduPilot's
Software-in-the-Loop (SITL) simulation. SITL runs the exact Pixhawk firmware on a
desktop PC and emulates the physical flight controller. The selected Pareto route is
uploaded as a waypoint file; SITL simulates the full flight and reports GPS track,
battery consumption profile, and route adherence.

\textbf{Pass criterion.} Route adherence within $\pm$2\,m of planned waypoints; mission
completion rate 100\%; mission upload time $<$30\,s end-to-end.

\subsection{Step 2: Hardware Integration on Raspberry Pi~5}

The complete Python pipeline (graph loader, Tier~1 pre-computation, Tier~2 selector,
MAVLink uploader) is deployed to the Raspberry Pi~5 mounted on the F450 quadcopter
frame. Integration tests verify that:
\begin{itemize}
  \item Tier~2 route selection and mission file generation complete in $<$1\,s.
  \item MAVLink mission upload to Pixhawk completes reliably.
  \item The Pixhawk acknowledges and stores the waypoint mission correctly.
\end{itemize}

\subsection{Step 3: Outdoor Flight Tests}

Real flights are conducted on the F450 quadcopter with a servo-actuated payload release
mechanism. Three flights are planned, one per payload preset. Each flight:
\begin{enumerate}
  \item Operator enters payload tags on the ground-station UI.
  \item Tier~2 selects route; mission file generated and uploaded via MAVLink.
  \item UAV executes the mission autonomously.
  \item GPS track is compared against the planned route.
  \item Payload is released at the delivery point.
\end{enumerate}

\textbf{Measured metrics.} End-to-end latency (tag entry to mission upload), route
adherence (average cross-track error), and payload delivery success.

\subsection{Summary}

Table~\ref{tab:demosummary} summarises the demonstration structure across both phases.

\begin{table}[ht!]
\centering
\caption{Demonstration Plan Summary}
\label{tab:demosummary}
\renewcommand{\arraystretch}{1.3}
\resizebox{\textwidth}{!}{%
\begin{tabular}{llll}
\toprule
\textbf{Phase} & \textbf{Demo} & \textbf{What is Shown} & \textbf{Environment} \\
\midrule
Phase I & 1 & Airspace graph \& objective visualisation & Simulation \\
Phase I & 2 & 3D Pareto front vs MOPSO baseline        & Simulation \\
Phase I & 3 & Interactive payload-to-route demo         & Dashboard / Laptop \\
Phase I & 4 & Quantitative benchmark table              & Simulation \\
Phase I & 5 & Tier~2 selection accuracy measurement     & Simulation \\
\midrule
Phase II & 6 & ArduPilot SITL end-to-end validation    & SITL / Desktop \\
Phase II & 7 & Raspberry Pi hardware integration test   & Onboard compute \\
Phase II & 8 & Outdoor flight tests (3 payload types)   & Real UAV outdoor \\
\bottomrule
\end{tabular}%
}
\end{table}
